% Save this file as chapter10.tex and add % Save this file as chapter10.tex and add % Save this file as chapter10.tex and add % Save this file as chapter10.tex and add \include{chapter10} in the main file.

\providecommand{\Sym}{\operatorname{Sym}}
\providecommand{\Core}{\operatorname{Core}}
\providecommand{\Fix}{\operatorname{Fix}}

\chapter{Permutation Actions and Coset Actions}
\label{chap:permutation-actions}


\section{Basic Notation for Permutation Actions}

\begin{definition}[Orbits and stabilizers]
Let \(G\leq \Sym(\Omega)\), where \(|\Omega|=n\).

\begin{enumerate}[label=\textup{(\arabic*)}]
    \item If \(\Delta\subseteq \Omega\) and \(g\in G\), define
    \[
        \Delta^g=\{\delta^g\mid \delta\in \Delta\}.
    \]

    \item For \(\omega_1,\omega_2\in \Omega\), write
    \[
        \omega_1\sim_G \omega_2
    \]
    if there exists \(g\in G\) such that \(\omega_1^g=\omega_2\).

    \item The orbit of \(\omega\in \Omega\) under \(G\) is
    \[
        \omega^G=\{\omega^g\mid g\in G\}.
    \]

    \item The stabilizer of \(\omega\in\Omega\) is
    \[
        G_\omega=\{g\in G\mid \omega^g=\omega\}.
    \]
\end{enumerate}
\end{definition}


\section{Setwise and Pointwise Stabilizers}

\begin{definition}[Stabilizers of subsets]
Let \(\Delta\subseteq \Omega\). Define the \emph{setwise stabilizer} of \(\Delta\) in \(G\) by
\[
    G_\Delta=\{g\in G\mid \Delta^g=\Delta\}.
\]
Define the \emph{pointwise stabilizer} of \(\Delta\) in \(G\) by
\[
    G_{(\Delta)}
    =
    \{g\in G\mid \delta^g=\delta \text{ for all } \delta\in\Delta\}.
\]
\end{definition}

\begin{proposition}
Let \(G\leq \Sym(\Omega)\) and let \(\Delta\subseteq \Omega\). Then
\[
    G_{(\Delta)}\unlhd G_\Delta.
\]
Moreover, \(G_\Delta\) acts on \(\Delta\), and the kernel of this induced action is \(G_{(\Delta)}\). Hence
\[
    G_\Delta^\Delta \iso G_\Delta/G_{(\Delta)},
\]
where \(G_\Delta^\Delta\) denotes the permutation group induced by \(G_\Delta\) on \(\Delta\).
\end{proposition}

\begin{proof}
First, \(G_{(\Delta)}\leq G_\Delta\), because every element fixing each point of \(\Delta\) certainly stabilizes \(\Delta\) as a set.

Let \(x\in G_\Delta\) and \(h\in G_{(\Delta)}\). We show that \(x^{-1}hx\in G_{(\Delta)}\). For any \(\delta\in \Delta\), since \(x\in G_\Delta\), we have \(\delta^{x^{-1}}\in \Delta\). Since \(h\) fixes every point of \(\Delta\),
\[
    (\delta^{x^{-1}})^h=\delta^{x^{-1}}.
\]
Therefore
\[
    \delta^{x^{-1}hx}
    =
    \bigl((\delta^{x^{-1}})^h\bigr)^x
    =
    (\delta^{x^{-1}})^x
    =
    \delta.
\]
Thus \(x^{-1}hx\in G_{(\Delta)}\), proving that
\[
    G_{(\Delta)}\unlhd G_\Delta.
\]

Now consider the restriction homomorphism
\[
    \rho:G_\Delta\longrightarrow \Sym(\Delta),
    \qquad
    g\longmapsto g|_\Delta.
\]
Its kernel consists precisely of the elements of \(G_\Delta\) that fix every point of \(\Delta\), namely \(G_{(\Delta)}\). Its image is \(G_\Delta^\Delta\). By the First Isomorphism Theorem,
\[
    G_\Delta/G_{(\Delta)}\iso G_\Delta^\Delta.
\]
\end{proof}

\begin{proposition}[Conjugating stabilizers]
Let \(G\leq \Sym(\Omega)\).

\begin{enumerate}[label=\textup{(\arabic*)}]
    \item For every \(\omega\in\Omega\) and \(g\in G\),
    \[
        g^{-1}G_\omega g=G_{\omega^g}.
    \]

    \item For every \(\Delta\subseteq\Omega\) and \(g\in G\),
    \[
        g^{-1}G_\Delta g=G_{\Delta^g}.
    \]

    \item For every \(\Delta\subseteq\Omega\) and \(g\in G\),
    \[
        g^{-1}G_{(\Delta)}g=G_{(\Delta^g)}.
    \]
\end{enumerate}
\end{proposition}

\begin{proof}
We prove the first statement. Let \(h\in G_\omega\). Then
\[
    (\omega^g)^{g^{-1}hg}
    =
    \omega^{hg}
    =
    (\omega^h)^g
    =
    \omega^g.
\]
Thus \(g^{-1}hg\in G_{\omega^g}\), so
\[
    g^{-1}G_\omega g\leq G_{\omega^g}.
\]
Applying the same argument with \(g^{-1}\) gives the reverse inclusion. Hence
\[
    g^{-1}G_\omega g=G_{\omega^g}.
\]

The proofs for \(G_\Delta\) and \(G_{(\Delta)}\) are identical, replacing the point \(\omega\) by the subset \(\Delta\).
\end{proof}

\section{Fixed Points and Semiregular Actions}

\begin{definition}[Fixed points]
Let \(G\) act on \(\Omega\).

\begin{enumerate}[label=\textup{(\arabic*)}]
    \item For \(g\in G\), define
    \[
        \Fix_\Omega(g)=\{\omega\in\Omega\mid \omega^g=\omega\}.
    \]

    \item The set of common fixed points of \(G\) on \(\Omega\) is
    \[
        \Fix_\Omega(G)
        =
        \{\omega\in\Omega\mid \omega^g=\omega \text{ for all } g\in G\}.
    \]
\end{enumerate}
\end{definition}

\begin{definition}[Semiregular and regular actions]
Let \(G\) act on \(\Omega\).

\begin{enumerate}[label=\textup{(\arabic*)}]
    \item The action is called \emph{semiregular} if
    \[
        G_\omega=1
    \]
    for every \(\omega\in\Omega\).

    \item The action is called \emph{regular} if it is both transitive and semiregular.
\end{enumerate}
\end{definition}

\begin{proposition}
Let \(G\) be finite and act semiregularly on \(\Omega\). Then every orbit has length \(|G|\). In particular, if \(\Omega\) is finite, then
\[
    |G|\mid |\Omega|.
\]
If the action is transitive and semiregular, then it is regular and
\[
    |G|=|\Omega|.
\]
\end{proposition}

\begin{proof}
Since the action is semiregular, \(G_\omega=1\) for every \(\omega\in\Omega\). By the orbit--stabilizer theorem,
\[
    |\omega^G|=|G:G_\omega|=|G|.
\]
Thus every orbit has length \(|G|\). If \(\Omega\) is finite, it is a disjoint union of orbits, so \(|G|\mid |\Omega|\).

If the action is transitive, there is only one orbit. Hence
\[
    |\Omega|=|\omega^G|=|G|.
\]
\end{proof}

\section{Coset Actions and the Core}

\begin{definition}[Right coset action]
Let \(G\) be a group and let \(H\leq G\). Let
\[
    [G:H]_{\mathrm r}:=\{Hg\mid g\in G\}
\]
be the set of right cosets of \(H\) in \(G\). Define an action of \(G\) on \([G:H]_{\mathrm r}\) by
\[
    (Hg)^x=Hgx
\]
for all \(g,x\in G\).
\end{definition}

\begin{proposition}
The right coset action of \(G\) on \([G:H]_{\mathrm r}\) is transitive.
\end{proposition}

\begin{proof}
For every coset \(Hg\in [G:H]_{\mathrm r}\),
\[
    H^g=Hg.
\]
Thus every coset lies in the orbit of \(H\), so the action is transitive.
\end{proof}

\begin{definition}[Core]
Let \(H\leq G\). The \emph{core} of \(H\) in \(G\) is
\[
    \Core_G(H)
    =
    \bigcap_{g\in G} g^{-1}Hg.
\]
\end{definition}

\begin{proposition}
Let \(G\) act on the set \([G:H]_{\mathrm r}\) by right multiplication. Then the kernel of this action is
\[
    \Core_G(H)=\bigcap_{g\in G}g^{-1}Hg.
\]
Moreover, \(\Core_G(H)\) is the largest normal subgroup of \(G\) contained in \(H\).
\end{proposition}

\begin{proof}
The stabilizer of the coset \(Hg\) is
\[
    G_{Hg}
    =
    \{x\in G\mid Hgx=Hg\}.
\]
Now
\[
    Hgx=Hg
    \iff
    gxg^{-1}\in H
    \iff
    x\in g^{-1}Hg.
\]
Thus
\[
    G_{Hg}=g^{-1}Hg.
\]
The kernel of the action is the intersection of all point stabilizers, so
\[
    \Ker(G\curvearrowright [G:H]_{\mathrm r})
    =
    \bigcap_{g\in G}G_{Hg}
    =
    \bigcap_{g\in G}g^{-1}Hg
    =
    \Core_G(H).
\]

Since the kernel of an action is normal in \(G\), \(\Core_G(H)\unlhd G\). Also, taking \(g=1\), we see that
\[
    \Core_G(H)\leq H.
\]

Finally, suppose \(N\unlhd G\) and \(N\leq H\). Then for every \(g\in G\),
\[
    N=g^{-1}Ng\leq g^{-1}Hg.
\]
Therefore
\[
    N\leq \bigcap_{g\in G}g^{-1}Hg=\Core_G(H).
\]
Thus \(\Core_G(H)\) is the largest normal subgroup of \(G\) contained in \(H\).
\end{proof}

\begin{proposition}
Let \(G\) act on
\[
    \Omega=[G:H]_{\mathrm r}
\]
by right multiplication. Then
\[
    G^\Omega \iso G/\Core_G(H).
\]
In particular, \(G^\Omega\) is a transitive permutation group of degree \(|G:H|\).
\end{proposition}

\begin{proof}
The action gives a homomorphism
\[
    \rho:G\longrightarrow \Sym(\Omega).
\]
Its image is \(G^\Omega\), and its kernel is \(\Core_G(H)\). Therefore, by the First Isomorphism Theorem,
\[
    G/\Core_G(H)\iso G^\Omega.
\]
The action is transitive by the previous proposition, and
\[
    |\Omega|=|G:H|.
\]
\end{proof}

\begin{definition}[Faithful action and core-free subgroup]
An action of \(G\) on \(\Omega\) is called \emph{faithful} if its kernel is trivial.

A subgroup \(H\leq G\) is called \emph{core-free} in \(G\) if
\[
    \Core_G(H)=1.
\]
Thus the right coset action of \(G\) on \([G:H]_{\mathrm r}\) is faithful if and only if \(H\) is core-free in \(G\).
\end{definition}

\section{Transitive Actions as Coset Actions}

\begin{definition}[Equivalence of actions]
Let \(G\) act on two sets \(\Omega\) and \(\Sigma\). The two actions are called \emph{equivalent} if there exists a bijection
\[
    \varphi:\Omega\longrightarrow \Sigma
\]
such that
\[
    (\omega^g)\varphi=(\omega\varphi)^g
\]
for every \(\omega\in\Omega\) and \(g\in G\).

Equivalently, the following diagram commutes for every \(g\in G\):
\[
\begin{tikzcd}
\omega \arrow[r, "\varphi"] \arrow[d, "{g}"'] &
\omega\varphi \arrow[d, "{g}"] \\
\omega^g \arrow[r, "\varphi"'] &
(\omega\varphi)^g.
\end{tikzcd}
\]
\end{definition}

\begin{theorem}[Every transitive action is a coset action]
Let \(G\leq \Sym(\Omega)\) be transitive. Fix \(\omega\in\Omega\), and put
\[
    H=G_\omega.
\]
Then the action of \(G\) on \(\Omega\) is equivalent to the right coset action of \(G\) on \([G:H]_{\mathrm r}\). Moreover, \(H\) is core-free in \(G\), so this coset action is faithful.
\end{theorem}

\begin{proof}
Let
\[
    \Sigma=[G:H]_{\mathrm r}.
\]
Define
\[
    \varphi:\Omega\longrightarrow \Sigma
\]
by
\[
    \varphi(\omega^g)=Hg.
\]
This is well-defined: if
\[
    \omega^g=\omega^h,
\]
then
\[
    \omega^{gh^{-1}}=\omega,
\]
so \(gh^{-1}\in G_\omega=H\), and hence
\[
    Hg=Hh.
\]

The map is surjective because every coset has the form \(Hg\), and it is injective by the same argument in reverse. Thus \(\varphi\) is a bijection.

For \(x\in G\),
\[
    \varphi\bigl((\omega^g)^x\bigr)
    =
    \varphi(\omega^{gx})
    =
    Hgx
    =
    (Hg)^x
    =
    \varphi(\omega^g)^x.
\]
Thus \(\varphi\) is an equivalence of actions.

It remains to show that \(H\) is core-free. Suppose \(N\unlhd G\) and \(N\leq H\). For every \(x\in G\),
\[
    N=N^x\leq H^x=G_{\omega^x}.
\]
Since \(G\) is transitive on \(\Omega\), every point of \(\Omega\) is of the form \(\omega^x\). Hence every element of \(N\) fixes every point of \(\Omega\). Since \(G\leq \Sym(\Omega)\), the action is faithful, so \(N=1\). Therefore
\[
    \Core_G(H)=1.
\]
\end{proof}

\begin{example}[Coset actions of \(S_4\)]
Let
\[
    G=S_4.
\]
Consider the following subgroups:
\[
    H_1=\gen{(12)},\qquad
    H_2=\gen{(12)(34)},\qquad
    H_3=\gen{(13)},\qquad
    H_4=\gen{(13)(24)}.
\]
For each \(i\), let
\[
    \Omega_i=[G:H_i]_{\mathrm r}.
\]

We compare the corresponding coset actions \(G^{\Omega_i}\).

\medskip

\noindent
\textbf{First, \(G^{\Omega_1}\) and \(G^{\Omega_2}\) are not equivalent.}

Suppose, for contradiction, that there exists an equivalence
\[
    \varphi:\Omega_1\longrightarrow \Omega_2.
\]
Write
\[
    H_1\varphi=H_2g
\]
for some \(g\in G\). Since \((12)\in H_1\), we have
\[
    H_1(12)=H_1.
\]
Using equivariance,
\[
    H_2g
    =
    H_1\varphi
    =
    (H_1(12))\varphi
    =
    (H_1\varphi)^{(12)}
    =
    (H_2g)^{(12)}
    =
    H_2g(12).
\]
Thus
\[
    H_2g=H_2g(12),
\]
so
\[
    g(12)g^{-1}\in H_2.
\]
But \(g(12)g^{-1}\) is a transposition, while
\[
    H_2=\{1,(12)(34)\}
\]
contains no transposition. This contradiction shows that \(G^{\Omega_1}\) and \(G^{\Omega_2}\) are not equivalent.

\medskip

\noindent
\textbf{Second, \(G^{\Omega_1}\) and \(G^{\Omega_3}\) are equivalent.}

Indeed,
\[
    H_3^{(23)}
    =
    (23)^{-1}H_3(23)
    =
    \gen{(12)}
    =
    H_1.
\]
Define
\[
    \varphi:\Omega_1\longrightarrow \Omega_3
\]
by
\[
    \varphi(H_1g)=H_3(23)g.
\]
We check that \(\varphi\) is well-defined. If
\[
    H_1g=H_1h,
\]
then \(gh^{-1}\in H_1\). Since
\[
    (23)H_1(23)=H_3,
\]
we get
\[
    (23)gh^{-1}(23)\in H_3.
\]
Hence
\[
    H_3(23)g=H_3(23)h.
\]
Thus \(\varphi\) is well-defined. It is plainly bijective, and for \(x\in G\),
\[
    \varphi\bigl((H_1g)^x\bigr)
    =
    \varphi(H_1gx)
    =
    H_3(23)gx
    =
    (H_3(23)g)^x
    =
    \varphi(H_1g)^x.
\]
Therefore the two actions are equivalent.

\medskip

Similarly, \(G^{\Omega_2}\) and \(G^{\Omega_4}\) are equivalent, since \(H_2\) and \(H_4\) are conjugate in \(S_4\). However, \(G^{\Omega_1}\) is not equivalent to \(G^{\Omega_2}\) or \(G^{\Omega_4}\), because a subgroup generated by a transposition cannot be conjugate to a subgroup generated by a double transposition.
\end{example}

\begin{theorem}[Criterion for equivalence of transitive actions]
Let \(G\) act transitively on two sets \(\Omega\) and \(\Sigma\). Fix \(\omega\in\Omega\) and \(\sigma\in\Sigma\). Then the two actions are equivalent if and only if the stabilizers \(G_\omega\) and \(G_\sigma\) are conjugate in \(G\).
\end{theorem}

\begin{proof}
Suppose first that the two actions are equivalent. Let
\[
    \varphi:\Omega\longrightarrow \Sigma
\]
be an equivalence. Since \(G\) is transitive on \(\Sigma\), there exists \(g\in G\) such that
\[
    \omega\varphi=\sigma^g.
\]
For \(x\in G_\omega\), we have \(\omega^x=\omega\). Hence
\[
    (\omega\varphi)^x
    =
    (\omega^x)\varphi
    =
    \omega\varphi.
\]
Thus \(x\in G_{\omega\varphi}\), so
\[
    G_\omega\leq G_{\omega\varphi}.
\]
Applying the same argument to \(\varphi^{-1}\), we get equality:
\[
    G_\omega=G_{\omega\varphi}.
\]
Since \(\omega\varphi=\sigma^g\), the conjugation formula for stabilizers gives
\[
    G_{\omega\varphi}
    =
    G_{\sigma^g}
    =
    g^{-1}G_\sigma g.
\]
Therefore
\[
    G_\omega=g^{-1}G_\sigma g,
\]
so \(G_\omega\) and \(G_\sigma\) are conjugate in \(G\).

Conversely, suppose that \(G_\omega\) and \(G_\sigma\) are conjugate. Choose \(g\in G\) such that
\[
    G_\omega=g^{-1}G_\sigma g.
\]
Define
\[
    \varphi:\Omega\longrightarrow \Sigma
\]
by
\[
    \varphi(\omega^x)=\sigma^{gx}.
\]

We check that this is well-defined. If
\[
    \omega^x=\omega^y,
\]
then
\[
    xy^{-1}\in G_\omega=g^{-1}G_\sigma g.
\]
Thus
\[
    gxy^{-1}g^{-1}\in G_\sigma,
\]
which implies
\[
    \sigma^{gx}=\sigma^{gy}.
\]
Hence \(\varphi\) is well-defined.

It is surjective because every element of \(\Sigma\) has the form \(\sigma^z\), and
\[
    \sigma^z=\sigma^{g(g^{-1}z)}=\varphi(\omega^{g^{-1}z}).
\]
It is injective by reversing the well-definedness argument.

Finally, for any \(a\in G\),
\[
    \varphi\bigl((\omega^x)^a\bigr)
    =
    \varphi(\omega^{xa})
    =
    \sigma^{gxa}
    =
    (\sigma^{gx})^a
    =
    \varphi(\omega^x)^a.
\]
Thus \(\varphi\) is an equivalence of actions.
\end{proof}

\begin{corollary}
The equivalence classes of transitive \(G\)-actions are in bijection with conjugacy classes of subgroups of \(G\).
\end{corollary}

\begin{proof}
Every transitive action is equivalent to the right coset action on \([G:H]_{\mathrm r}\), where \(H\) is a point stabilizer. By the previous theorem, two such coset actions are equivalent exactly when their stabilizers are conjugate in \(G\).
\end{proof} in the main file.

\providecommand{\Sym}{\operatorname{Sym}}
\providecommand{\Core}{\operatorname{Core}}
\providecommand{\Fix}{\operatorname{Fix}}

\chapter{Permutation Actions and Coset Actions}
\label{chap:permutation-actions}


\section{Basic Notation for Permutation Actions}

\begin{definition}[Orbits and stabilizers]
Let \(G\leq \Sym(\Omega)\), where \(|\Omega|=n\).

\begin{enumerate}[label=\textup{(\arabic*)}]
    \item If \(\Delta\subseteq \Omega\) and \(g\in G\), define
    \[
        \Delta^g=\{\delta^g\mid \delta\in \Delta\}.
    \]

    \item For \(\omega_1,\omega_2\in \Omega\), write
    \[
        \omega_1\sim_G \omega_2
    \]
    if there exists \(g\in G\) such that \(\omega_1^g=\omega_2\).

    \item The orbit of \(\omega\in \Omega\) under \(G\) is
    \[
        \omega^G=\{\omega^g\mid g\in G\}.
    \]

    \item The stabilizer of \(\omega\in\Omega\) is
    \[
        G_\omega=\{g\in G\mid \omega^g=\omega\}.
    \]
\end{enumerate}
\end{definition}


\section{Setwise and Pointwise Stabilizers}

\begin{definition}[Stabilizers of subsets]
Let \(\Delta\subseteq \Omega\). Define the \emph{setwise stabilizer} of \(\Delta\) in \(G\) by
\[
    G_\Delta=\{g\in G\mid \Delta^g=\Delta\}.
\]
Define the \emph{pointwise stabilizer} of \(\Delta\) in \(G\) by
\[
    G_{(\Delta)}
    =
    \{g\in G\mid \delta^g=\delta \text{ for all } \delta\in\Delta\}.
\]
\end{definition}

\begin{proposition}
Let \(G\leq \Sym(\Omega)\) and let \(\Delta\subseteq \Omega\). Then
\[
    G_{(\Delta)}\unlhd G_\Delta.
\]
Moreover, \(G_\Delta\) acts on \(\Delta\), and the kernel of this induced action is \(G_{(\Delta)}\). Hence
\[
    G_\Delta^\Delta \iso G_\Delta/G_{(\Delta)},
\]
where \(G_\Delta^\Delta\) denotes the permutation group induced by \(G_\Delta\) on \(\Delta\).
\end{proposition}

\begin{proof}
First, \(G_{(\Delta)}\leq G_\Delta\), because every element fixing each point of \(\Delta\) certainly stabilizes \(\Delta\) as a set.

Let \(x\in G_\Delta\) and \(h\in G_{(\Delta)}\). We show that \(x^{-1}hx\in G_{(\Delta)}\). For any \(\delta\in \Delta\), since \(x\in G_\Delta\), we have \(\delta^{x^{-1}}\in \Delta\). Since \(h\) fixes every point of \(\Delta\),
\[
    (\delta^{x^{-1}})^h=\delta^{x^{-1}}.
\]
Therefore
\[
    \delta^{x^{-1}hx}
    =
    \bigl((\delta^{x^{-1}})^h\bigr)^x
    =
    (\delta^{x^{-1}})^x
    =
    \delta.
\]
Thus \(x^{-1}hx\in G_{(\Delta)}\), proving that
\[
    G_{(\Delta)}\unlhd G_\Delta.
\]

Now consider the restriction homomorphism
\[
    \rho:G_\Delta\longrightarrow \Sym(\Delta),
    \qquad
    g\longmapsto g|_\Delta.
\]
Its kernel consists precisely of the elements of \(G_\Delta\) that fix every point of \(\Delta\), namely \(G_{(\Delta)}\). Its image is \(G_\Delta^\Delta\). By the First Isomorphism Theorem,
\[
    G_\Delta/G_{(\Delta)}\iso G_\Delta^\Delta.
\]
\end{proof}

\begin{proposition}[Conjugating stabilizers]
Let \(G\leq \Sym(\Omega)\).

\begin{enumerate}[label=\textup{(\arabic*)}]
    \item For every \(\omega\in\Omega\) and \(g\in G\),
    \[
        g^{-1}G_\omega g=G_{\omega^g}.
    \]

    \item For every \(\Delta\subseteq\Omega\) and \(g\in G\),
    \[
        g^{-1}G_\Delta g=G_{\Delta^g}.
    \]

    \item For every \(\Delta\subseteq\Omega\) and \(g\in G\),
    \[
        g^{-1}G_{(\Delta)}g=G_{(\Delta^g)}.
    \]
\end{enumerate}
\end{proposition}

\begin{proof}
We prove the first statement. Let \(h\in G_\omega\). Then
\[
    (\omega^g)^{g^{-1}hg}
    =
    \omega^{hg}
    =
    (\omega^h)^g
    =
    \omega^g.
\]
Thus \(g^{-1}hg\in G_{\omega^g}\), so
\[
    g^{-1}G_\omega g\leq G_{\omega^g}.
\]
Applying the same argument with \(g^{-1}\) gives the reverse inclusion. Hence
\[
    g^{-1}G_\omega g=G_{\omega^g}.
\]

The proofs for \(G_\Delta\) and \(G_{(\Delta)}\) are identical, replacing the point \(\omega\) by the subset \(\Delta\).
\end{proof}

\section{Fixed Points and Semiregular Actions}

\begin{definition}[Fixed points]
Let \(G\) act on \(\Omega\).

\begin{enumerate}[label=\textup{(\arabic*)}]
    \item For \(g\in G\), define
    \[
        \Fix_\Omega(g)=\{\omega\in\Omega\mid \omega^g=\omega\}.
    \]

    \item The set of common fixed points of \(G\) on \(\Omega\) is
    \[
        \Fix_\Omega(G)
        =
        \{\omega\in\Omega\mid \omega^g=\omega \text{ for all } g\in G\}.
    \]
\end{enumerate}
\end{definition}

\begin{definition}[Semiregular and regular actions]
Let \(G\) act on \(\Omega\).

\begin{enumerate}[label=\textup{(\arabic*)}]
    \item The action is called \emph{semiregular} if
    \[
        G_\omega=1
    \]
    for every \(\omega\in\Omega\).

    \item The action is called \emph{regular} if it is both transitive and semiregular.
\end{enumerate}
\end{definition}

\begin{proposition}
Let \(G\) be finite and act semiregularly on \(\Omega\). Then every orbit has length \(|G|\). In particular, if \(\Omega\) is finite, then
\[
    |G|\mid |\Omega|.
\]
If the action is transitive and semiregular, then it is regular and
\[
    |G|=|\Omega|.
\]
\end{proposition}

\begin{proof}
Since the action is semiregular, \(G_\omega=1\) for every \(\omega\in\Omega\). By the orbit--stabilizer theorem,
\[
    |\omega^G|=|G:G_\omega|=|G|.
\]
Thus every orbit has length \(|G|\). If \(\Omega\) is finite, it is a disjoint union of orbits, so \(|G|\mid |\Omega|\).

If the action is transitive, there is only one orbit. Hence
\[
    |\Omega|=|\omega^G|=|G|.
\]
\end{proof}

\section{Coset Actions and the Core}

\begin{definition}[Right coset action]
Let \(G\) be a group and let \(H\leq G\). Let
\[
    [G:H]_{\mathrm r}:=\{Hg\mid g\in G\}
\]
be the set of right cosets of \(H\) in \(G\). Define an action of \(G\) on \([G:H]_{\mathrm r}\) by
\[
    (Hg)^x=Hgx
\]
for all \(g,x\in G\).
\end{definition}

\begin{proposition}
The right coset action of \(G\) on \([G:H]_{\mathrm r}\) is transitive.
\end{proposition}

\begin{proof}
For every coset \(Hg\in [G:H]_{\mathrm r}\),
\[
    H^g=Hg.
\]
Thus every coset lies in the orbit of \(H\), so the action is transitive.
\end{proof}

\begin{definition}[Core]
Let \(H\leq G\). The \emph{core} of \(H\) in \(G\) is
\[
    \Core_G(H)
    =
    \bigcap_{g\in G} g^{-1}Hg.
\]
\end{definition}

\begin{proposition}
Let \(G\) act on the set \([G:H]_{\mathrm r}\) by right multiplication. Then the kernel of this action is
\[
    \Core_G(H)=\bigcap_{g\in G}g^{-1}Hg.
\]
Moreover, \(\Core_G(H)\) is the largest normal subgroup of \(G\) contained in \(H\).
\end{proposition}

\begin{proof}
The stabilizer of the coset \(Hg\) is
\[
    G_{Hg}
    =
    \{x\in G\mid Hgx=Hg\}.
\]
Now
\[
    Hgx=Hg
    \iff
    gxg^{-1}\in H
    \iff
    x\in g^{-1}Hg.
\]
Thus
\[
    G_{Hg}=g^{-1}Hg.
\]
The kernel of the action is the intersection of all point stabilizers, so
\[
    \Ker(G\curvearrowright [G:H]_{\mathrm r})
    =
    \bigcap_{g\in G}G_{Hg}
    =
    \bigcap_{g\in G}g^{-1}Hg
    =
    \Core_G(H).
\]

Since the kernel of an action is normal in \(G\), \(\Core_G(H)\unlhd G\). Also, taking \(g=1\), we see that
\[
    \Core_G(H)\leq H.
\]

Finally, suppose \(N\unlhd G\) and \(N\leq H\). Then for every \(g\in G\),
\[
    N=g^{-1}Ng\leq g^{-1}Hg.
\]
Therefore
\[
    N\leq \bigcap_{g\in G}g^{-1}Hg=\Core_G(H).
\]
Thus \(\Core_G(H)\) is the largest normal subgroup of \(G\) contained in \(H\).
\end{proof}

\begin{proposition}
Let \(G\) act on
\[
    \Omega=[G:H]_{\mathrm r}
\]
by right multiplication. Then
\[
    G^\Omega \iso G/\Core_G(H).
\]
In particular, \(G^\Omega\) is a transitive permutation group of degree \(|G:H|\).
\end{proposition}

\begin{proof}
The action gives a homomorphism
\[
    \rho:G\longrightarrow \Sym(\Omega).
\]
Its image is \(G^\Omega\), and its kernel is \(\Core_G(H)\). Therefore, by the First Isomorphism Theorem,
\[
    G/\Core_G(H)\iso G^\Omega.
\]
The action is transitive by the previous proposition, and
\[
    |\Omega|=|G:H|.
\]
\end{proof}

\begin{definition}[Faithful action and core-free subgroup]
An action of \(G\) on \(\Omega\) is called \emph{faithful} if its kernel is trivial.

A subgroup \(H\leq G\) is called \emph{core-free} in \(G\) if
\[
    \Core_G(H)=1.
\]
Thus the right coset action of \(G\) on \([G:H]_{\mathrm r}\) is faithful if and only if \(H\) is core-free in \(G\).
\end{definition}

\section{Transitive Actions as Coset Actions}

\begin{definition}[Equivalence of actions]
Let \(G\) act on two sets \(\Omega\) and \(\Sigma\). The two actions are called \emph{equivalent} if there exists a bijection
\[
    \varphi:\Omega\longrightarrow \Sigma
\]
such that
\[
    (\omega^g)\varphi=(\omega\varphi)^g
\]
for every \(\omega\in\Omega\) and \(g\in G\).

Equivalently, the following diagram commutes for every \(g\in G\):
\[
\begin{tikzcd}
\omega \arrow[r, "\varphi"] \arrow[d, "{g}"'] &
\omega\varphi \arrow[d, "{g}"] \\
\omega^g \arrow[r, "\varphi"'] &
(\omega\varphi)^g.
\end{tikzcd}
\]
\end{definition}

\begin{theorem}[Every transitive action is a coset action]
Let \(G\leq \Sym(\Omega)\) be transitive. Fix \(\omega\in\Omega\), and put
\[
    H=G_\omega.
\]
Then the action of \(G\) on \(\Omega\) is equivalent to the right coset action of \(G\) on \([G:H]_{\mathrm r}\). Moreover, \(H\) is core-free in \(G\), so this coset action is faithful.
\end{theorem}

\begin{proof}
Let
\[
    \Sigma=[G:H]_{\mathrm r}.
\]
Define
\[
    \varphi:\Omega\longrightarrow \Sigma
\]
by
\[
    \varphi(\omega^g)=Hg.
\]
This is well-defined: if
\[
    \omega^g=\omega^h,
\]
then
\[
    \omega^{gh^{-1}}=\omega,
\]
so \(gh^{-1}\in G_\omega=H\), and hence
\[
    Hg=Hh.
\]

The map is surjective because every coset has the form \(Hg\), and it is injective by the same argument in reverse. Thus \(\varphi\) is a bijection.

For \(x\in G\),
\[
    \varphi\bigl((\omega^g)^x\bigr)
    =
    \varphi(\omega^{gx})
    =
    Hgx
    =
    (Hg)^x
    =
    \varphi(\omega^g)^x.
\]
Thus \(\varphi\) is an equivalence of actions.

It remains to show that \(H\) is core-free. Suppose \(N\unlhd G\) and \(N\leq H\). For every \(x\in G\),
\[
    N=N^x\leq H^x=G_{\omega^x}.
\]
Since \(G\) is transitive on \(\Omega\), every point of \(\Omega\) is of the form \(\omega^x\). Hence every element of \(N\) fixes every point of \(\Omega\). Since \(G\leq \Sym(\Omega)\), the action is faithful, so \(N=1\). Therefore
\[
    \Core_G(H)=1.
\]
\end{proof}

\begin{example}[Coset actions of \(S_4\)]
Let
\[
    G=S_4.
\]
Consider the following subgroups:
\[
    H_1=\gen{(12)},\qquad
    H_2=\gen{(12)(34)},\qquad
    H_3=\gen{(13)},\qquad
    H_4=\gen{(13)(24)}.
\]
For each \(i\), let
\[
    \Omega_i=[G:H_i]_{\mathrm r}.
\]

We compare the corresponding coset actions \(G^{\Omega_i}\).

\medskip

\noindent
\textbf{First, \(G^{\Omega_1}\) and \(G^{\Omega_2}\) are not equivalent.}

Suppose, for contradiction, that there exists an equivalence
\[
    \varphi:\Omega_1\longrightarrow \Omega_2.
\]
Write
\[
    H_1\varphi=H_2g
\]
for some \(g\in G\). Since \((12)\in H_1\), we have
\[
    H_1(12)=H_1.
\]
Using equivariance,
\[
    H_2g
    =
    H_1\varphi
    =
    (H_1(12))\varphi
    =
    (H_1\varphi)^{(12)}
    =
    (H_2g)^{(12)}
    =
    H_2g(12).
\]
Thus
\[
    H_2g=H_2g(12),
\]
so
\[
    g(12)g^{-1}\in H_2.
\]
But \(g(12)g^{-1}\) is a transposition, while
\[
    H_2=\{1,(12)(34)\}
\]
contains no transposition. This contradiction shows that \(G^{\Omega_1}\) and \(G^{\Omega_2}\) are not equivalent.

\medskip

\noindent
\textbf{Second, \(G^{\Omega_1}\) and \(G^{\Omega_3}\) are equivalent.}

Indeed,
\[
    H_3^{(23)}
    =
    (23)^{-1}H_3(23)
    =
    \gen{(12)}
    =
    H_1.
\]
Define
\[
    \varphi:\Omega_1\longrightarrow \Omega_3
\]
by
\[
    \varphi(H_1g)=H_3(23)g.
\]
We check that \(\varphi\) is well-defined. If
\[
    H_1g=H_1h,
\]
then \(gh^{-1}\in H_1\). Since
\[
    (23)H_1(23)=H_3,
\]
we get
\[
    (23)gh^{-1}(23)\in H_3.
\]
Hence
\[
    H_3(23)g=H_3(23)h.
\]
Thus \(\varphi\) is well-defined. It is plainly bijective, and for \(x\in G\),
\[
    \varphi\bigl((H_1g)^x\bigr)
    =
    \varphi(H_1gx)
    =
    H_3(23)gx
    =
    (H_3(23)g)^x
    =
    \varphi(H_1g)^x.
\]
Therefore the two actions are equivalent.

\medskip

Similarly, \(G^{\Omega_2}\) and \(G^{\Omega_4}\) are equivalent, since \(H_2\) and \(H_4\) are conjugate in \(S_4\). However, \(G^{\Omega_1}\) is not equivalent to \(G^{\Omega_2}\) or \(G^{\Omega_4}\), because a subgroup generated by a transposition cannot be conjugate to a subgroup generated by a double transposition.
\end{example}

\begin{theorem}[Criterion for equivalence of transitive actions]
Let \(G\) act transitively on two sets \(\Omega\) and \(\Sigma\). Fix \(\omega\in\Omega\) and \(\sigma\in\Sigma\). Then the two actions are equivalent if and only if the stabilizers \(G_\omega\) and \(G_\sigma\) are conjugate in \(G\).
\end{theorem}

\begin{proof}
Suppose first that the two actions are equivalent. Let
\[
    \varphi:\Omega\longrightarrow \Sigma
\]
be an equivalence. Since \(G\) is transitive on \(\Sigma\), there exists \(g\in G\) such that
\[
    \omega\varphi=\sigma^g.
\]
For \(x\in G_\omega\), we have \(\omega^x=\omega\). Hence
\[
    (\omega\varphi)^x
    =
    (\omega^x)\varphi
    =
    \omega\varphi.
\]
Thus \(x\in G_{\omega\varphi}\), so
\[
    G_\omega\leq G_{\omega\varphi}.
\]
Applying the same argument to \(\varphi^{-1}\), we get equality:
\[
    G_\omega=G_{\omega\varphi}.
\]
Since \(\omega\varphi=\sigma^g\), the conjugation formula for stabilizers gives
\[
    G_{\omega\varphi}
    =
    G_{\sigma^g}
    =
    g^{-1}G_\sigma g.
\]
Therefore
\[
    G_\omega=g^{-1}G_\sigma g,
\]
so \(G_\omega\) and \(G_\sigma\) are conjugate in \(G\).

Conversely, suppose that \(G_\omega\) and \(G_\sigma\) are conjugate. Choose \(g\in G\) such that
\[
    G_\omega=g^{-1}G_\sigma g.
\]
Define
\[
    \varphi:\Omega\longrightarrow \Sigma
\]
by
\[
    \varphi(\omega^x)=\sigma^{gx}.
\]

We check that this is well-defined. If
\[
    \omega^x=\omega^y,
\]
then
\[
    xy^{-1}\in G_\omega=g^{-1}G_\sigma g.
\]
Thus
\[
    gxy^{-1}g^{-1}\in G_\sigma,
\]
which implies
\[
    \sigma^{gx}=\sigma^{gy}.
\]
Hence \(\varphi\) is well-defined.

It is surjective because every element of \(\Sigma\) has the form \(\sigma^z\), and
\[
    \sigma^z=\sigma^{g(g^{-1}z)}=\varphi(\omega^{g^{-1}z}).
\]
It is injective by reversing the well-definedness argument.

Finally, for any \(a\in G\),
\[
    \varphi\bigl((\omega^x)^a\bigr)
    =
    \varphi(\omega^{xa})
    =
    \sigma^{gxa}
    =
    (\sigma^{gx})^a
    =
    \varphi(\omega^x)^a.
\]
Thus \(\varphi\) is an equivalence of actions.
\end{proof}

\begin{corollary}
The equivalence classes of transitive \(G\)-actions are in bijection with conjugacy classes of subgroups of \(G\).
\end{corollary}

\begin{proof}
Every transitive action is equivalent to the right coset action on \([G:H]_{\mathrm r}\), where \(H\) is a point stabilizer. By the previous theorem, two such coset actions are equivalent exactly when their stabilizers are conjugate in \(G\).
\end{proof} in the main file.

\providecommand{\Sym}{\operatorname{Sym}}
\providecommand{\Core}{\operatorname{Core}}
\providecommand{\Fix}{\operatorname{Fix}}

\chapter{Permutation Actions and Coset Actions}
\label{chap:permutation-actions}


\section{Basic Notation for Permutation Actions}

\begin{definition}[Orbits and stabilizers]
Let \(G\leq \Sym(\Omega)\), where \(|\Omega|=n\).

\begin{enumerate}[label=\textup{(\arabic*)}]
    \item If \(\Delta\subseteq \Omega\) and \(g\in G\), define
    \[
        \Delta^g=\{\delta^g\mid \delta\in \Delta\}.
    \]

    \item For \(\omega_1,\omega_2\in \Omega\), write
    \[
        \omega_1\sim_G \omega_2
    \]
    if there exists \(g\in G\) such that \(\omega_1^g=\omega_2\).

    \item The orbit of \(\omega\in \Omega\) under \(G\) is
    \[
        \omega^G=\{\omega^g\mid g\in G\}.
    \]

    \item The stabilizer of \(\omega\in\Omega\) is
    \[
        G_\omega=\{g\in G\mid \omega^g=\omega\}.
    \]
\end{enumerate}
\end{definition}


\section{Setwise and Pointwise Stabilizers}

\begin{definition}[Stabilizers of subsets]
Let \(\Delta\subseteq \Omega\). Define the \emph{setwise stabilizer} of \(\Delta\) in \(G\) by
\[
    G_\Delta=\{g\in G\mid \Delta^g=\Delta\}.
\]
Define the \emph{pointwise stabilizer} of \(\Delta\) in \(G\) by
\[
    G_{(\Delta)}
    =
    \{g\in G\mid \delta^g=\delta \text{ for all } \delta\in\Delta\}.
\]
\end{definition}

\begin{proposition}
Let \(G\leq \Sym(\Omega)\) and let \(\Delta\subseteq \Omega\). Then
\[
    G_{(\Delta)}\unlhd G_\Delta.
\]
Moreover, \(G_\Delta\) acts on \(\Delta\), and the kernel of this induced action is \(G_{(\Delta)}\). Hence
\[
    G_\Delta^\Delta \iso G_\Delta/G_{(\Delta)},
\]
where \(G_\Delta^\Delta\) denotes the permutation group induced by \(G_\Delta\) on \(\Delta\).
\end{proposition}

\begin{proof}
First, \(G_{(\Delta)}\leq G_\Delta\), because every element fixing each point of \(\Delta\) certainly stabilizes \(\Delta\) as a set.

Let \(x\in G_\Delta\) and \(h\in G_{(\Delta)}\). We show that \(x^{-1}hx\in G_{(\Delta)}\). For any \(\delta\in \Delta\), since \(x\in G_\Delta\), we have \(\delta^{x^{-1}}\in \Delta\). Since \(h\) fixes every point of \(\Delta\),
\[
    (\delta^{x^{-1}})^h=\delta^{x^{-1}}.
\]
Therefore
\[
    \delta^{x^{-1}hx}
    =
    \bigl((\delta^{x^{-1}})^h\bigr)^x
    =
    (\delta^{x^{-1}})^x
    =
    \delta.
\]
Thus \(x^{-1}hx\in G_{(\Delta)}\), proving that
\[
    G_{(\Delta)}\unlhd G_\Delta.
\]

Now consider the restriction homomorphism
\[
    \rho:G_\Delta\longrightarrow \Sym(\Delta),
    \qquad
    g\longmapsto g|_\Delta.
\]
Its kernel consists precisely of the elements of \(G_\Delta\) that fix every point of \(\Delta\), namely \(G_{(\Delta)}\). Its image is \(G_\Delta^\Delta\). By the First Isomorphism Theorem,
\[
    G_\Delta/G_{(\Delta)}\iso G_\Delta^\Delta.
\]
\end{proof}

\begin{proposition}[Conjugating stabilizers]
Let \(G\leq \Sym(\Omega)\).

\begin{enumerate}[label=\textup{(\arabic*)}]
    \item For every \(\omega\in\Omega\) and \(g\in G\),
    \[
        g^{-1}G_\omega g=G_{\omega^g}.
    \]

    \item For every \(\Delta\subseteq\Omega\) and \(g\in G\),
    \[
        g^{-1}G_\Delta g=G_{\Delta^g}.
    \]

    \item For every \(\Delta\subseteq\Omega\) and \(g\in G\),
    \[
        g^{-1}G_{(\Delta)}g=G_{(\Delta^g)}.
    \]
\end{enumerate}
\end{proposition}

\begin{proof}
We prove the first statement. Let \(h\in G_\omega\). Then
\[
    (\omega^g)^{g^{-1}hg}
    =
    \omega^{hg}
    =
    (\omega^h)^g
    =
    \omega^g.
\]
Thus \(g^{-1}hg\in G_{\omega^g}\), so
\[
    g^{-1}G_\omega g\leq G_{\omega^g}.
\]
Applying the same argument with \(g^{-1}\) gives the reverse inclusion. Hence
\[
    g^{-1}G_\omega g=G_{\omega^g}.
\]

The proofs for \(G_\Delta\) and \(G_{(\Delta)}\) are identical, replacing the point \(\omega\) by the subset \(\Delta\).
\end{proof}

\section{Fixed Points and Semiregular Actions}

\begin{definition}[Fixed points]
Let \(G\) act on \(\Omega\).

\begin{enumerate}[label=\textup{(\arabic*)}]
    \item For \(g\in G\), define
    \[
        \Fix_\Omega(g)=\{\omega\in\Omega\mid \omega^g=\omega\}.
    \]

    \item The set of common fixed points of \(G\) on \(\Omega\) is
    \[
        \Fix_\Omega(G)
        =
        \{\omega\in\Omega\mid \omega^g=\omega \text{ for all } g\in G\}.
    \]
\end{enumerate}
\end{definition}

\begin{definition}[Semiregular and regular actions]
Let \(G\) act on \(\Omega\).

\begin{enumerate}[label=\textup{(\arabic*)}]
    \item The action is called \emph{semiregular} if
    \[
        G_\omega=1
    \]
    for every \(\omega\in\Omega\).

    \item The action is called \emph{regular} if it is both transitive and semiregular.
\end{enumerate}
\end{definition}

\begin{proposition}
Let \(G\) be finite and act semiregularly on \(\Omega\). Then every orbit has length \(|G|\). In particular, if \(\Omega\) is finite, then
\[
    |G|\mid |\Omega|.
\]
If the action is transitive and semiregular, then it is regular and
\[
    |G|=|\Omega|.
\]
\end{proposition}

\begin{proof}
Since the action is semiregular, \(G_\omega=1\) for every \(\omega\in\Omega\). By the orbit--stabilizer theorem,
\[
    |\omega^G|=|G:G_\omega|=|G|.
\]
Thus every orbit has length \(|G|\). If \(\Omega\) is finite, it is a disjoint union of orbits, so \(|G|\mid |\Omega|\).

If the action is transitive, there is only one orbit. Hence
\[
    |\Omega|=|\omega^G|=|G|.
\]
\end{proof}

\section{Coset Actions and the Core}

\begin{definition}[Right coset action]
Let \(G\) be a group and let \(H\leq G\). Let
\[
    [G:H]_{\mathrm r}:=\{Hg\mid g\in G\}
\]
be the set of right cosets of \(H\) in \(G\). Define an action of \(G\) on \([G:H]_{\mathrm r}\) by
\[
    (Hg)^x=Hgx
\]
for all \(g,x\in G\).
\end{definition}

\begin{proposition}
The right coset action of \(G\) on \([G:H]_{\mathrm r}\) is transitive.
\end{proposition}

\begin{proof}
For every coset \(Hg\in [G:H]_{\mathrm r}\),
\[
    H^g=Hg.
\]
Thus every coset lies in the orbit of \(H\), so the action is transitive.
\end{proof}

\begin{definition}[Core]
Let \(H\leq G\). The \emph{core} of \(H\) in \(G\) is
\[
    \Core_G(H)
    =
    \bigcap_{g\in G} g^{-1}Hg.
\]
\end{definition}

\begin{proposition}
Let \(G\) act on the set \([G:H]_{\mathrm r}\) by right multiplication. Then the kernel of this action is
\[
    \Core_G(H)=\bigcap_{g\in G}g^{-1}Hg.
\]
Moreover, \(\Core_G(H)\) is the largest normal subgroup of \(G\) contained in \(H\).
\end{proposition}

\begin{proof}
The stabilizer of the coset \(Hg\) is
\[
    G_{Hg}
    =
    \{x\in G\mid Hgx=Hg\}.
\]
Now
\[
    Hgx=Hg
    \iff
    gxg^{-1}\in H
    \iff
    x\in g^{-1}Hg.
\]
Thus
\[
    G_{Hg}=g^{-1}Hg.
\]
The kernel of the action is the intersection of all point stabilizers, so
\[
    \Ker(G\curvearrowright [G:H]_{\mathrm r})
    =
    \bigcap_{g\in G}G_{Hg}
    =
    \bigcap_{g\in G}g^{-1}Hg
    =
    \Core_G(H).
\]

Since the kernel of an action is normal in \(G\), \(\Core_G(H)\unlhd G\). Also, taking \(g=1\), we see that
\[
    \Core_G(H)\leq H.
\]

Finally, suppose \(N\unlhd G\) and \(N\leq H\). Then for every \(g\in G\),
\[
    N=g^{-1}Ng\leq g^{-1}Hg.
\]
Therefore
\[
    N\leq \bigcap_{g\in G}g^{-1}Hg=\Core_G(H).
\]
Thus \(\Core_G(H)\) is the largest normal subgroup of \(G\) contained in \(H\).
\end{proof}

\begin{proposition}
Let \(G\) act on
\[
    \Omega=[G:H]_{\mathrm r}
\]
by right multiplication. Then
\[
    G^\Omega \iso G/\Core_G(H).
\]
In particular, \(G^\Omega\) is a transitive permutation group of degree \(|G:H|\).
\end{proposition}

\begin{proof}
The action gives a homomorphism
\[
    \rho:G\longrightarrow \Sym(\Omega).
\]
Its image is \(G^\Omega\), and its kernel is \(\Core_G(H)\). Therefore, by the First Isomorphism Theorem,
\[
    G/\Core_G(H)\iso G^\Omega.
\]
The action is transitive by the previous proposition, and
\[
    |\Omega|=|G:H|.
\]
\end{proof}

\begin{definition}[Faithful action and core-free subgroup]
An action of \(G\) on \(\Omega\) is called \emph{faithful} if its kernel is trivial.

A subgroup \(H\leq G\) is called \emph{core-free} in \(G\) if
\[
    \Core_G(H)=1.
\]
Thus the right coset action of \(G\) on \([G:H]_{\mathrm r}\) is faithful if and only if \(H\) is core-free in \(G\).
\end{definition}

\section{Transitive Actions as Coset Actions}

\begin{definition}[Equivalence of actions]
Let \(G\) act on two sets \(\Omega\) and \(\Sigma\). The two actions are called \emph{equivalent} if there exists a bijection
\[
    \varphi:\Omega\longrightarrow \Sigma
\]
such that
\[
    (\omega^g)\varphi=(\omega\varphi)^g
\]
for every \(\omega\in\Omega\) and \(g\in G\).

Equivalently, the following diagram commutes for every \(g\in G\):
\[
\begin{tikzcd}
\omega \arrow[r, "\varphi"] \arrow[d, "{g}"'] &
\omega\varphi \arrow[d, "{g}"] \\
\omega^g \arrow[r, "\varphi"'] &
(\omega\varphi)^g.
\end{tikzcd}
\]
\end{definition}

\begin{theorem}[Every transitive action is a coset action]
Let \(G\leq \Sym(\Omega)\) be transitive. Fix \(\omega\in\Omega\), and put
\[
    H=G_\omega.
\]
Then the action of \(G\) on \(\Omega\) is equivalent to the right coset action of \(G\) on \([G:H]_{\mathrm r}\). Moreover, \(H\) is core-free in \(G\), so this coset action is faithful.
\end{theorem}

\begin{proof}
Let
\[
    \Sigma=[G:H]_{\mathrm r}.
\]
Define
\[
    \varphi:\Omega\longrightarrow \Sigma
\]
by
\[
    \varphi(\omega^g)=Hg.
\]
This is well-defined: if
\[
    \omega^g=\omega^h,
\]
then
\[
    \omega^{gh^{-1}}=\omega,
\]
so \(gh^{-1}\in G_\omega=H\), and hence
\[
    Hg=Hh.
\]

The map is surjective because every coset has the form \(Hg\), and it is injective by the same argument in reverse. Thus \(\varphi\) is a bijection.

For \(x\in G\),
\[
    \varphi\bigl((\omega^g)^x\bigr)
    =
    \varphi(\omega^{gx})
    =
    Hgx
    =
    (Hg)^x
    =
    \varphi(\omega^g)^x.
\]
Thus \(\varphi\) is an equivalence of actions.

It remains to show that \(H\) is core-free. Suppose \(N\unlhd G\) and \(N\leq H\). For every \(x\in G\),
\[
    N=N^x\leq H^x=G_{\omega^x}.
\]
Since \(G\) is transitive on \(\Omega\), every point of \(\Omega\) is of the form \(\omega^x\). Hence every element of \(N\) fixes every point of \(\Omega\). Since \(G\leq \Sym(\Omega)\), the action is faithful, so \(N=1\). Therefore
\[
    \Core_G(H)=1.
\]
\end{proof}

\begin{example}[Coset actions of \(S_4\)]
Let
\[
    G=S_4.
\]
Consider the following subgroups:
\[
    H_1=\gen{(12)},\qquad
    H_2=\gen{(12)(34)},\qquad
    H_3=\gen{(13)},\qquad
    H_4=\gen{(13)(24)}.
\]
For each \(i\), let
\[
    \Omega_i=[G:H_i]_{\mathrm r}.
\]

We compare the corresponding coset actions \(G^{\Omega_i}\).

\medskip

\noindent
\textbf{First, \(G^{\Omega_1}\) and \(G^{\Omega_2}\) are not equivalent.}

Suppose, for contradiction, that there exists an equivalence
\[
    \varphi:\Omega_1\longrightarrow \Omega_2.
\]
Write
\[
    H_1\varphi=H_2g
\]
for some \(g\in G\). Since \((12)\in H_1\), we have
\[
    H_1(12)=H_1.
\]
Using equivariance,
\[
    H_2g
    =
    H_1\varphi
    =
    (H_1(12))\varphi
    =
    (H_1\varphi)^{(12)}
    =
    (H_2g)^{(12)}
    =
    H_2g(12).
\]
Thus
\[
    H_2g=H_2g(12),
\]
so
\[
    g(12)g^{-1}\in H_2.
\]
But \(g(12)g^{-1}\) is a transposition, while
\[
    H_2=\{1,(12)(34)\}
\]
contains no transposition. This contradiction shows that \(G^{\Omega_1}\) and \(G^{\Omega_2}\) are not equivalent.

\medskip

\noindent
\textbf{Second, \(G^{\Omega_1}\) and \(G^{\Omega_3}\) are equivalent.}

Indeed,
\[
    H_3^{(23)}
    =
    (23)^{-1}H_3(23)
    =
    \gen{(12)}
    =
    H_1.
\]
Define
\[
    \varphi:\Omega_1\longrightarrow \Omega_3
\]
by
\[
    \varphi(H_1g)=H_3(23)g.
\]
We check that \(\varphi\) is well-defined. If
\[
    H_1g=H_1h,
\]
then \(gh^{-1}\in H_1\). Since
\[
    (23)H_1(23)=H_3,
\]
we get
\[
    (23)gh^{-1}(23)\in H_3.
\]
Hence
\[
    H_3(23)g=H_3(23)h.
\]
Thus \(\varphi\) is well-defined. It is plainly bijective, and for \(x\in G\),
\[
    \varphi\bigl((H_1g)^x\bigr)
    =
    \varphi(H_1gx)
    =
    H_3(23)gx
    =
    (H_3(23)g)^x
    =
    \varphi(H_1g)^x.
\]
Therefore the two actions are equivalent.

\medskip

Similarly, \(G^{\Omega_2}\) and \(G^{\Omega_4}\) are equivalent, since \(H_2\) and \(H_4\) are conjugate in \(S_4\). However, \(G^{\Omega_1}\) is not equivalent to \(G^{\Omega_2}\) or \(G^{\Omega_4}\), because a subgroup generated by a transposition cannot be conjugate to a subgroup generated by a double transposition.
\end{example}

\begin{theorem}[Criterion for equivalence of transitive actions]
Let \(G\) act transitively on two sets \(\Omega\) and \(\Sigma\). Fix \(\omega\in\Omega\) and \(\sigma\in\Sigma\). Then the two actions are equivalent if and only if the stabilizers \(G_\omega\) and \(G_\sigma\) are conjugate in \(G\).
\end{theorem}

\begin{proof}
Suppose first that the two actions are equivalent. Let
\[
    \varphi:\Omega\longrightarrow \Sigma
\]
be an equivalence. Since \(G\) is transitive on \(\Sigma\), there exists \(g\in G\) such that
\[
    \omega\varphi=\sigma^g.
\]
For \(x\in G_\omega\), we have \(\omega^x=\omega\). Hence
\[
    (\omega\varphi)^x
    =
    (\omega^x)\varphi
    =
    \omega\varphi.
\]
Thus \(x\in G_{\omega\varphi}\), so
\[
    G_\omega\leq G_{\omega\varphi}.
\]
Applying the same argument to \(\varphi^{-1}\), we get equality:
\[
    G_\omega=G_{\omega\varphi}.
\]
Since \(\omega\varphi=\sigma^g\), the conjugation formula for stabilizers gives
\[
    G_{\omega\varphi}
    =
    G_{\sigma^g}
    =
    g^{-1}G_\sigma g.
\]
Therefore
\[
    G_\omega=g^{-1}G_\sigma g,
\]
so \(G_\omega\) and \(G_\sigma\) are conjugate in \(G\).

Conversely, suppose that \(G_\omega\) and \(G_\sigma\) are conjugate. Choose \(g\in G\) such that
\[
    G_\omega=g^{-1}G_\sigma g.
\]
Define
\[
    \varphi:\Omega\longrightarrow \Sigma
\]
by
\[
    \varphi(\omega^x)=\sigma^{gx}.
\]

We check that this is well-defined. If
\[
    \omega^x=\omega^y,
\]
then
\[
    xy^{-1}\in G_\omega=g^{-1}G_\sigma g.
\]
Thus
\[
    gxy^{-1}g^{-1}\in G_\sigma,
\]
which implies
\[
    \sigma^{gx}=\sigma^{gy}.
\]
Hence \(\varphi\) is well-defined.

It is surjective because every element of \(\Sigma\) has the form \(\sigma^z\), and
\[
    \sigma^z=\sigma^{g(g^{-1}z)}=\varphi(\omega^{g^{-1}z}).
\]
It is injective by reversing the well-definedness argument.

Finally, for any \(a\in G\),
\[
    \varphi\bigl((\omega^x)^a\bigr)
    =
    \varphi(\omega^{xa})
    =
    \sigma^{gxa}
    =
    (\sigma^{gx})^a
    =
    \varphi(\omega^x)^a.
\]
Thus \(\varphi\) is an equivalence of actions.
\end{proof}

\begin{corollary}
The equivalence classes of transitive \(G\)-actions are in bijection with conjugacy classes of subgroups of \(G\).
\end{corollary}

\begin{proof}
Every transitive action is equivalent to the right coset action on \([G:H]_{\mathrm r}\), where \(H\) is a point stabilizer. By the previous theorem, two such coset actions are equivalent exactly when their stabilizers are conjugate in \(G\).
\end{proof} in the main file.

\providecommand{\Sym}{\operatorname{Sym}}
\providecommand{\Core}{\operatorname{Core}}
\providecommand{\Fix}{\operatorname{Fix}}

\chapter{Permutation Actions and Coset Actions}
\label{chap:permutation-actions}


\section{Basic Notation for Permutation Actions}

\begin{definition}[Orbits and stabilizers]
Let \(G\leq \Sym(\Omega)\), where \(|\Omega|=n\).

\begin{enumerate}[label=\textup{(\arabic*)}]
    \item If \(\Delta\subseteq \Omega\) and \(g\in G\), define
    \[
        \Delta^g=\{\delta^g\mid \delta\in \Delta\}.
    \]

    \item For \(\omega_1,\omega_2\in \Omega\), write
    \[
        \omega_1\sim_G \omega_2
    \]
    if there exists \(g\in G\) such that \(\omega_1^g=\omega_2\).

    \item The orbit of \(\omega\in \Omega\) under \(G\) is
    \[
        \omega^G=\{\omega^g\mid g\in G\}.
    \]

    \item The stabilizer of \(\omega\in\Omega\) is
    \[
        G_\omega=\{g\in G\mid \omega^g=\omega\}.
    \]
\end{enumerate}
\end{definition}


\section{Setwise and Pointwise Stabilizers}

\begin{definition}[Stabilizers of subsets]
Let \(\Delta\subseteq \Omega\). Define the \emph{setwise stabilizer} of \(\Delta\) in \(G\) by
\[
    G_\Delta=\{g\in G\mid \Delta^g=\Delta\}.
\]
Define the \emph{pointwise stabilizer} of \(\Delta\) in \(G\) by
\[
    G_{(\Delta)}
    =
    \{g\in G\mid \delta^g=\delta \text{ for all } \delta\in\Delta\}.
\]
\end{definition}

\begin{proposition}
Let \(G\leq \Sym(\Omega)\) and let \(\Delta\subseteq \Omega\). Then
\[
    G_{(\Delta)}\unlhd G_\Delta.
\]
Moreover, \(G_\Delta\) acts on \(\Delta\), and the kernel of this induced action is \(G_{(\Delta)}\). Hence
\[
    G_\Delta^\Delta \iso G_\Delta/G_{(\Delta)},
\]
where \(G_\Delta^\Delta\) denotes the permutation group induced by \(G_\Delta\) on \(\Delta\).
\end{proposition}

\begin{proof}
First, \(G_{(\Delta)}\leq G_\Delta\), because every element fixing each point of \(\Delta\) certainly stabilizes \(\Delta\) as a set.

Let \(x\in G_\Delta\) and \(h\in G_{(\Delta)}\). We show that \(x^{-1}hx\in G_{(\Delta)}\). For any \(\delta\in \Delta\), since \(x\in G_\Delta\), we have \(\delta^{x^{-1}}\in \Delta\). Since \(h\) fixes every point of \(\Delta\),
\[
    (\delta^{x^{-1}})^h=\delta^{x^{-1}}.
\]
Therefore
\[
    \delta^{x^{-1}hx}
    =
    \bigl((\delta^{x^{-1}})^h\bigr)^x
    =
    (\delta^{x^{-1}})^x
    =
    \delta.
\]
Thus \(x^{-1}hx\in G_{(\Delta)}\), proving that
\[
    G_{(\Delta)}\unlhd G_\Delta.
\]

Now consider the restriction homomorphism
\[
    \rho:G_\Delta\longrightarrow \Sym(\Delta),
    \qquad
    g\longmapsto g|_\Delta.
\]
Its kernel consists precisely of the elements of \(G_\Delta\) that fix every point of \(\Delta\), namely \(G_{(\Delta)}\). Its image is \(G_\Delta^\Delta\). By the First Isomorphism Theorem,
\[
    G_\Delta/G_{(\Delta)}\iso G_\Delta^\Delta.
\]
\end{proof}

\begin{proposition}[Conjugating stabilizers]
Let \(G\leq \Sym(\Omega)\).

\begin{enumerate}[label=\textup{(\arabic*)}]
    \item For every \(\omega\in\Omega\) and \(g\in G\),
    \[
        g^{-1}G_\omega g=G_{\omega^g}.
    \]

    \item For every \(\Delta\subseteq\Omega\) and \(g\in G\),
    \[
        g^{-1}G_\Delta g=G_{\Delta^g}.
    \]

    \item For every \(\Delta\subseteq\Omega\) and \(g\in G\),
    \[
        g^{-1}G_{(\Delta)}g=G_{(\Delta^g)}.
    \]
\end{enumerate}
\end{proposition}

\begin{proof}
We prove the first statement. Let \(h\in G_\omega\). Then
\[
    (\omega^g)^{g^{-1}hg}
    =
    \omega^{hg}
    =
    (\omega^h)^g
    =
    \omega^g.
\]
Thus \(g^{-1}hg\in G_{\omega^g}\), so
\[
    g^{-1}G_\omega g\leq G_{\omega^g}.
\]
Applying the same argument with \(g^{-1}\) gives the reverse inclusion. Hence
\[
    g^{-1}G_\omega g=G_{\omega^g}.
\]

The proofs for \(G_\Delta\) and \(G_{(\Delta)}\) are identical, replacing the point \(\omega\) by the subset \(\Delta\).
\end{proof}

\section{Fixed Points and Semiregular Actions}

\begin{definition}[Fixed points]
Let \(G\) act on \(\Omega\).

\begin{enumerate}[label=\textup{(\arabic*)}]
    \item For \(g\in G\), define
    \[
        \Fix_\Omega(g)=\{\omega\in\Omega\mid \omega^g=\omega\}.
    \]

    \item The set of common fixed points of \(G\) on \(\Omega\) is
    \[
        \Fix_\Omega(G)
        =
        \{\omega\in\Omega\mid \omega^g=\omega \text{ for all } g\in G\}.
    \]
\end{enumerate}
\end{definition}

\begin{definition}[Semiregular and regular actions]
Let \(G\) act on \(\Omega\).

\begin{enumerate}[label=\textup{(\arabic*)}]
    \item The action is called \emph{semiregular} if
    \[
        G_\omega=1
    \]
    for every \(\omega\in\Omega\).

    \item The action is called \emph{regular} if it is both transitive and semiregular.
\end{enumerate}
\end{definition}

\begin{proposition}
Let \(G\) be finite and act semiregularly on \(\Omega\). Then every orbit has length \(|G|\). In particular, if \(\Omega\) is finite, then
\[
    |G|\mid |\Omega|.
\]
If the action is transitive and semiregular, then it is regular and
\[
    |G|=|\Omega|.
\]
\end{proposition}

\begin{proof}
Since the action is semiregular, \(G_\omega=1\) for every \(\omega\in\Omega\). By the orbit--stabilizer theorem,
\[
    |\omega^G|=|G:G_\omega|=|G|.
\]
Thus every orbit has length \(|G|\). If \(\Omega\) is finite, it is a disjoint union of orbits, so \(|G|\mid |\Omega|\).

If the action is transitive, there is only one orbit. Hence
\[
    |\Omega|=|\omega^G|=|G|.
\]
\end{proof}

\section{Coset Actions and the Core}

\begin{definition}[Right coset action]
Let \(G\) be a group and let \(H\leq G\). Let
\[
    [G:H]_{\mathrm r}:=\{Hg\mid g\in G\}
\]
be the set of right cosets of \(H\) in \(G\). Define an action of \(G\) on \([G:H]_{\mathrm r}\) by
\[
    (Hg)^x=Hgx
\]
for all \(g,x\in G\).
\end{definition}

\begin{proposition}
The right coset action of \(G\) on \([G:H]_{\mathrm r}\) is transitive.
\end{proposition}

\begin{proof}
For every coset \(Hg\in [G:H]_{\mathrm r}\),
\[
    H^g=Hg.
\]
Thus every coset lies in the orbit of \(H\), so the action is transitive.
\end{proof}

\begin{definition}[Core]
Let \(H\leq G\). The \emph{core} of \(H\) in \(G\) is
\[
    \Core_G(H)
    =
    \bigcap_{g\in G} g^{-1}Hg.
\]
\end{definition}

\begin{proposition}
Let \(G\) act on the set \([G:H]_{\mathrm r}\) by right multiplication. Then the kernel of this action is
\[
    \Core_G(H)=\bigcap_{g\in G}g^{-1}Hg.
\]
Moreover, \(\Core_G(H)\) is the largest normal subgroup of \(G\) contained in \(H\).
\end{proposition}

\begin{proof}
The stabilizer of the coset \(Hg\) is
\[
    G_{Hg}
    =
    \{x\in G\mid Hgx=Hg\}.
\]
Now
\[
    Hgx=Hg
    \iff
    gxg^{-1}\in H
    \iff
    x\in g^{-1}Hg.
\]
Thus
\[
    G_{Hg}=g^{-1}Hg.
\]
The kernel of the action is the intersection of all point stabilizers, so
\[
    \Ker(G\curvearrowright [G:H]_{\mathrm r})
    =
    \bigcap_{g\in G}G_{Hg}
    =
    \bigcap_{g\in G}g^{-1}Hg
    =
    \Core_G(H).
\]

Since the kernel of an action is normal in \(G\), \(\Core_G(H)\unlhd G\). Also, taking \(g=1\), we see that
\[
    \Core_G(H)\leq H.
\]

Finally, suppose \(N\unlhd G\) and \(N\leq H\). Then for every \(g\in G\),
\[
    N=g^{-1}Ng\leq g^{-1}Hg.
\]
Therefore
\[
    N\leq \bigcap_{g\in G}g^{-1}Hg=\Core_G(H).
\]
Thus \(\Core_G(H)\) is the largest normal subgroup of \(G\) contained in \(H\).
\end{proof}

\begin{proposition}
Let \(G\) act on
\[
    \Omega=[G:H]_{\mathrm r}
\]
by right multiplication. Then
\[
    G^\Omega \iso G/\Core_G(H).
\]
In particular, \(G^\Omega\) is a transitive permutation group of degree \(|G:H|\).
\end{proposition}

\begin{proof}
The action gives a homomorphism
\[
    \rho:G\longrightarrow \Sym(\Omega).
\]
Its image is \(G^\Omega\), and its kernel is \(\Core_G(H)\). Therefore, by the First Isomorphism Theorem,
\[
    G/\Core_G(H)\iso G^\Omega.
\]
The action is transitive by the previous proposition, and
\[
    |\Omega|=|G:H|.
\]
\end{proof}

\begin{definition}[Faithful action and core-free subgroup]
An action of \(G\) on \(\Omega\) is called \emph{faithful} if its kernel is trivial.

A subgroup \(H\leq G\) is called \emph{core-free} in \(G\) if
\[
    \Core_G(H)=1.
\]
Thus the right coset action of \(G\) on \([G:H]_{\mathrm r}\) is faithful if and only if \(H\) is core-free in \(G\).
\end{definition}

\section{Transitive Actions as Coset Actions}

\begin{definition}[Equivalence of actions]
Let \(G\) act on two sets \(\Omega\) and \(\Sigma\). The two actions are called \emph{equivalent} if there exists a bijection
\[
    \varphi:\Omega\longrightarrow \Sigma
\]
such that
\[
    (\omega^g)\varphi=(\omega\varphi)^g
\]
for every \(\omega\in\Omega\) and \(g\in G\).

Equivalently, the following diagram commutes for every \(g\in G\):
\[
\begin{tikzcd}
\omega \arrow[r, "\varphi"] \arrow[d, "{g}"'] &
\omega\varphi \arrow[d, "{g}"] \\
\omega^g \arrow[r, "\varphi"'] &
(\omega\varphi)^g.
\end{tikzcd}
\]
\end{definition}

\begin{theorem}[Every transitive action is a coset action]
Let \(G\leq \Sym(\Omega)\) be transitive. Fix \(\omega\in\Omega\), and put
\[
    H=G_\omega.
\]
Then the action of \(G\) on \(\Omega\) is equivalent to the right coset action of \(G\) on \([G:H]_{\mathrm r}\). Moreover, \(H\) is core-free in \(G\), so this coset action is faithful.
\end{theorem}

\begin{proof}
Let
\[
    \Sigma=[G:H]_{\mathrm r}.
\]
Define
\[
    \varphi:\Omega\longrightarrow \Sigma
\]
by
\[
    \varphi(\omega^g)=Hg.
\]
This is well-defined: if
\[
    \omega^g=\omega^h,
\]
then
\[
    \omega^{gh^{-1}}=\omega,
\]
so \(gh^{-1}\in G_\omega=H\), and hence
\[
    Hg=Hh.
\]

The map is surjective because every coset has the form \(Hg\), and it is injective by the same argument in reverse. Thus \(\varphi\) is a bijection.

For \(x\in G\),
\[
    \varphi\bigl((\omega^g)^x\bigr)
    =
    \varphi(\omega^{gx})
    =
    Hgx
    =
    (Hg)^x
    =
    \varphi(\omega^g)^x.
\]
Thus \(\varphi\) is an equivalence of actions.

It remains to show that \(H\) is core-free. Suppose \(N\unlhd G\) and \(N\leq H\). For every \(x\in G\),
\[
    N=N^x\leq H^x=G_{\omega^x}.
\]
Since \(G\) is transitive on \(\Omega\), every point of \(\Omega\) is of the form \(\omega^x\). Hence every element of \(N\) fixes every point of \(\Omega\). Since \(G\leq \Sym(\Omega)\), the action is faithful, so \(N=1\). Therefore
\[
    \Core_G(H)=1.
\]
\end{proof}

\begin{example}[Coset actions of \(S_4\)]
Let
\[
    G=S_4.
\]
Consider the following subgroups:
\[
    H_1=\gen{(12)},\qquad
    H_2=\gen{(12)(34)},\qquad
    H_3=\gen{(13)},\qquad
    H_4=\gen{(13)(24)}.
\]
For each \(i\), let
\[
    \Omega_i=[G:H_i]_{\mathrm r}.
\]

We compare the corresponding coset actions \(G^{\Omega_i}\).

\medskip

\noindent
\textbf{First, \(G^{\Omega_1}\) and \(G^{\Omega_2}\) are not equivalent.}

Suppose, for contradiction, that there exists an equivalence
\[
    \varphi:\Omega_1\longrightarrow \Omega_2.
\]
Write
\[
    H_1\varphi=H_2g
\]
for some \(g\in G\). Since \((12)\in H_1\), we have
\[
    H_1(12)=H_1.
\]
Using equivariance,
\[
    H_2g
    =
    H_1\varphi
    =
    (H_1(12))\varphi
    =
    (H_1\varphi)^{(12)}
    =
    (H_2g)^{(12)}
    =
    H_2g(12).
\]
Thus
\[
    H_2g=H_2g(12),
\]
so
\[
    g(12)g^{-1}\in H_2.
\]
But \(g(12)g^{-1}\) is a transposition, while
\[
    H_2=\{1,(12)(34)\}
\]
contains no transposition. This contradiction shows that \(G^{\Omega_1}\) and \(G^{\Omega_2}\) are not equivalent.

\medskip

\noindent
\textbf{Second, \(G^{\Omega_1}\) and \(G^{\Omega_3}\) are equivalent.}

Indeed,
\[
    H_3^{(23)}
    =
    (23)^{-1}H_3(23)
    =
    \gen{(12)}
    =
    H_1.
\]
Define
\[
    \varphi:\Omega_1\longrightarrow \Omega_3
\]
by
\[
    \varphi(H_1g)=H_3(23)g.
\]
We check that \(\varphi\) is well-defined. If
\[
    H_1g=H_1h,
\]
then \(gh^{-1}\in H_1\). Since
\[
    (23)H_1(23)=H_3,
\]
we get
\[
    (23)gh^{-1}(23)\in H_3.
\]
Hence
\[
    H_3(23)g=H_3(23)h.
\]
Thus \(\varphi\) is well-defined. It is plainly bijective, and for \(x\in G\),
\[
    \varphi\bigl((H_1g)^x\bigr)
    =
    \varphi(H_1gx)
    =
    H_3(23)gx
    =
    (H_3(23)g)^x
    =
    \varphi(H_1g)^x.
\]
Therefore the two actions are equivalent.

\medskip

Similarly, \(G^{\Omega_2}\) and \(G^{\Omega_4}\) are equivalent, since \(H_2\) and \(H_4\) are conjugate in \(S_4\). However, \(G^{\Omega_1}\) is not equivalent to \(G^{\Omega_2}\) or \(G^{\Omega_4}\), because a subgroup generated by a transposition cannot be conjugate to a subgroup generated by a double transposition.
\end{example}

\begin{theorem}[Criterion for equivalence of transitive actions]
Let \(G\) act transitively on two sets \(\Omega\) and \(\Sigma\). Fix \(\omega\in\Omega\) and \(\sigma\in\Sigma\). Then the two actions are equivalent if and only if the stabilizers \(G_\omega\) and \(G_\sigma\) are conjugate in \(G\).
\end{theorem}

\begin{proof}
Suppose first that the two actions are equivalent. Let
\[
    \varphi:\Omega\longrightarrow \Sigma
\]
be an equivalence. Since \(G\) is transitive on \(\Sigma\), there exists \(g\in G\) such that
\[
    \omega\varphi=\sigma^g.
\]
For \(x\in G_\omega\), we have \(\omega^x=\omega\). Hence
\[
    (\omega\varphi)^x
    =
    (\omega^x)\varphi
    =
    \omega\varphi.
\]
Thus \(x\in G_{\omega\varphi}\), so
\[
    G_\omega\leq G_{\omega\varphi}.
\]
Applying the same argument to \(\varphi^{-1}\), we get equality:
\[
    G_\omega=G_{\omega\varphi}.
\]
Since \(\omega\varphi=\sigma^g\), the conjugation formula for stabilizers gives
\[
    G_{\omega\varphi}
    =
    G_{\sigma^g}
    =
    g^{-1}G_\sigma g.
\]
Therefore
\[
    G_\omega=g^{-1}G_\sigma g,
\]
so \(G_\omega\) and \(G_\sigma\) are conjugate in \(G\).

Conversely, suppose that \(G_\omega\) and \(G_\sigma\) are conjugate. Choose \(g\in G\) such that
\[
    G_\omega=g^{-1}G_\sigma g.
\]
Define
\[
    \varphi:\Omega\longrightarrow \Sigma
\]
by
\[
    \varphi(\omega^x)=\sigma^{gx}.
\]

We check that this is well-defined. If
\[
    \omega^x=\omega^y,
\]
then
\[
    xy^{-1}\in G_\omega=g^{-1}G_\sigma g.
\]
Thus
\[
    gxy^{-1}g^{-1}\in G_\sigma,
\]
which implies
\[
    \sigma^{gx}=\sigma^{gy}.
\]
Hence \(\varphi\) is well-defined.

It is surjective because every element of \(\Sigma\) has the form \(\sigma^z\), and
\[
    \sigma^z=\sigma^{g(g^{-1}z)}=\varphi(\omega^{g^{-1}z}).
\]
It is injective by reversing the well-definedness argument.

Finally, for any \(a\in G\),
\[
    \varphi\bigl((\omega^x)^a\bigr)
    =
    \varphi(\omega^{xa})
    =
    \sigma^{gxa}
    =
    (\sigma^{gx})^a
    =
    \varphi(\omega^x)^a.
\]
Thus \(\varphi\) is an equivalence of actions.
\end{proof}

\begin{corollary}
The equivalence classes of transitive \(G\)-actions are in bijection with conjugacy classes of subgroups of \(G\).
\end{corollary}

\begin{proof}
Every transitive action is equivalent to the right coset action on \([G:H]_{\mathrm r}\), where \(H\) is a point stabilizer. By the previous theorem, two such coset actions are equivalent exactly when their stabilizers are conjugate in \(G\).
\end{proof}